% =============================================================================
% Semana 10 — Trabalho 5 + Tableau Interativo: Filtros, Parâmetros e Ações
% Aula de 3 horas (19h–22h) | 06/05/2026
% Observação: 06/05 é feriado (Vitória de Santo Antão) — aula somente segunda (04/05)
% =============================================================================

\chapter{Semana 10 (06/05/2026) — Tableau Interativo: Filtros, Parâmetros e Ações}

% ---------------------------------------------------------------------------
\section{Objetivo da semana}
% ---------------------------------------------------------------------------

Pessoal, até aqui construímos visuais estáticos com narrativa clara.
Hoje aprendemos como torná-los \textbf{interativos} — dashboards que
respondem às perguntas do usuário em tempo real, sem precisar criar uma
versão diferente para cada recorte.

Filtros, parâmetros e ações são os mecanismos que transformam um dashboard
de relatório para \textbf{ferramenta de exploração}.

\begin{NoteBox}
\textbf{Calendário:} 06/05 é feriado de Vitória de Santo Antão.
Aula concentrada na segunda-feira (04/05).
\end{NoteBox}

\begin{BoardBox}
\textbf{Roteiro da aula (19h–22h)}
\begin{enumerate}
  \item 19h–19h40 — \textbf{Apresentações:} Trabalho 5 (storytelling + Tableau)
  \item 19h40–20h10 — \textbf{TEORIA:} Filtros no Tableau — tipos e hierarquia
  \item 20h10–20h40 — \textbf{TEORIA:} Parâmetros — criação e uso em campos calculados
  \item 20h40–21h — \textbf{TEORIA:} Ações de Dashboard — filter, highlight, URL, navigate
  \item 21h–21h40 — \textbf{PRÁTICA GUIADA:} Dashboard interativo com World Happiness Report
  \item 21h40–21h50 — \textbf{ESTUDO ESTATÍSTICO:} Correlação PIB × felicidade + outliers
  \item 21h50–22h — \textbf{Lançamento do T6 + revisão Q3}
\end{enumerate}
\end{BoardBox}

% =============================================================================
\section{Dataset da Semana — World Happiness Report (Kaggle)}
% =============================================================================

\begin{BoardBox}
\textbf{Base de dados da semana:} \textit{World Happiness Report}

\textbf{Link para download:}
\url{https://www.kaggle.com/datasets/unsdsn/world-happiness}

\textbf{Sobre o dataset:}
\begin{itemize}
  \item $\approx$ 150 países com índice de felicidade e fatores contribuintes
  \item \textbf{Colunas principais:} Country, Region, Happiness Score,
  GDP per Capita, Social Support, Healthy Life Expectancy, Freedom,
  Generosity, Perceptions of Corruption
  \item Dados limpos, sem joins necessários — perfeito para mapas e filtros
  \item Componente geográfico: ideal para mapa coroplético no Tableau
\end{itemize}

\textbf{Pergunta de negócio da semana:}
\textit{``Quais fatores mais impactam a felicidade de um país?
 PIB é suficiente ou outros fatores são mais relevantes?''}
\end{BoardBox}

\begin{NoteBox}
\textbf{Importação no Tableau:} File $\to$ Open $\to$ Text file $\to$
\texttt{2019.csv} (ou o ano mais recente). Verificar se \texttt{Country}
é reconhecido como Geographic Role (Country/Region). Se não:
clique direito em Country $\to$ Geographic Role $\to$ Country.
\end{NoteBox}

% =============================================================================
\section{Parte I — Apresentações: Trabalho 5 (19h–19h40)}
% =============================================================================

\begin{BoardBox}
\textbf{Checklist de avaliação T5 — Storytelling com Dados:}
\begin{tabular}{|p{4.5cm}|p{8.5cm}|}
\hline
\textbf{Ponto} & \textbf{O que observar} \\
\hline
Big Idea presente & 1 frase: ponto de vista + tensão + stake \\
Pergunta de Decisão & Explicitada no início da apresentação \\
Storyboard executado & Visual segue a sequência planejada \\
Títulos-mensagem & Nenhum título genérico (ex.: ``Gráfico 1'') \\
Hierarquia visual & 1 elemento de destaque por tela; cor com propósito \\
Tableau funcional & Pelo menos 2 visuais interativos \\
Pitch 3 min & Apresentação dentro do tempo; call to action ao final \\
\hline
\end{tabular}
\end{BoardBox}

% =============================================================================
\section{Parte II — Filtros no Tableau (19h40–20h10)}
% =============================================================================

\begin{BoardBox}
\textbf{Hierarquia de filtros no Tableau (para o quadro)}

O Tableau executa filtros em uma ordem específica. Entender essa ordem
evita resultados inesperados:

\begin{enumerate}
  \item \textbf{Extract filter:} aplicado na extração dos dados — filtra antes de carregar.
  \item \textbf{Data source filter:} filtra no nível da fonte de dados.
  \item \textbf{Context filter:} define um subconjunto de dados que outros filtros
  vão usar como base. Essencial para filtros dependentes (Top N dentro de uma região).
  \item \textbf{Dimension filter:} filtra dimensões (categorias).
  \item \textbf{Measure filter / Table calculation filter:} filtros sobre métricas
  e cálculos de tabela.
\end{enumerate}

\textbf{Tipo mais comum de confusão:} Top N sem Context Filter.
O Tableau calcula o Top N sobre o dataset inteiro, mesmo que você já tenha
aplicado um filtro de região. A solução: adicionar o filtro de região como
\emph{Context Filter} antes.
\end{BoardBox}

\begin{SolvedBox}
\textbf{Tipos de filtros para o usuário final (Quick Filters):}

\begin{tabular}{|l|p{9cm}|}
\hline
\textbf{Tipo} & \textbf{Quando usar} \\
\hline
Lista (checkbox) & Categoria com poucos valores; seleção múltipla \\
Lista suspensa & Categoria com muitos valores; seleção única \\
Controle deslizante & Variável numérica (ex.: ano, faixa de valor) \\
Data relativa & Datas (hoje, últimos 7 dias, mês corrente) \\
Pesquisa & Categoria com muitos valores únicos (busca por texto) \\
\hline
\end{tabular}

\textbf{Boas práticas:}
\begin{itemize}
  \item Nunca exponha mais de 4–5 filtros simultâneos para o usuário — excesso paralisa.
  \item Dê nomes claros: ``Período de Análise'' em vez de ``Order Date''.
  \item Defina um valor padrão que já mostre um recorte útil ao abrir o dashboard.
\end{itemize}
\end{SolvedBox}

% =============================================================================
\section{Parte III — Parâmetros (20h10–20h40)}
% =============================================================================

\begin{BoardBox}
\textbf{O que é um parâmetro no Tableau? (para o quadro)}

Um parâmetro é um \textbf{valor dinâmico controlado pelo usuário} que pode ser
usado em campos calculados, filtros e eixos de referência. Diferente de um filtro
(que inclui/exclui linhas), um parâmetro \textbf{muda o cálculo} do campo.

\textbf{Exemplos de uso:}
\begin{itemize}
  \item Parâmetro de ``Top N'': usuário escolhe ver Top 5, Top 10 ou Top 20 produtos.
  \item Parâmetro de meta: usuário insere a meta de vendas e o dashboard mostra quais
  regiões estão acima ou abaixo dinamicamente.
  \item Parâmetro de granularidade: alternar entre visão por dia, semana ou mês.
  \item Parâmetro de métrica: usuário escolhe visualizar por Receita ou por Quantidade.
\end{itemize}
\end{BoardBox}

\begin{SolvedBox}
\textbf{Como criar um parâmetro básico (passo a passo):}

\begin{enumerate}
  \item No painel de dados, clique com o botão direito $\to$ ``Criar Parâmetro''.
  \item Defina: tipo (inteiro/float/string/data), faixa de valores ou lista.
  \item Clique em ``Exibir Controle'' para mostrar o seletor ao usuário.
  \item No campo calculado: use \texttt{[Nome do Parâmetro]} como um valor dinâmico.
\end{enumerate}

\textbf{Exemplo de campo calculado com parâmetro de meta:}\\
\texttt{IF [Vendas] >= [Meta Vendas] THEN "Atingiu" ELSE "Abaixo" END}

O campo \texttt{[Meta Vendas]} é um parâmetro. O usuário ajusta a meta em
tempo real e o gráfico de desempenho atualiza imediatamente.
\end{SolvedBox}

% =============================================================================
\section{LOD Expressions — Introdução (complemento de Parâmetros)}
% =============================================================================

\begin{BoardBox}
\textbf{O que são LOD Expressions? (Level of Detail)}

No Tableau, a granularidade do cálculo normalmente depende das dimensões
na view. As \textbf{LOD expressions} permitem calcular métricas em um
nível de detalhe \textbf{diferente} do que está na tela.

\textbf{Os 3 tipos:}
\begin{tabular}{|l|p{5cm}|p{5.5cm}|}
\hline
\textbf{LOD} & \textbf{O que faz} & \textbf{Exemplo} \\
\hline
\texttt{FIXED} & Calcula no nível especificado,
ignorando filtros da view &
\texttt{\{FIXED [Department] : AVG([Salary])\}} —
média salarial por departamento, independente de filtros \\
\hline
\texttt{INCLUDE} & Adiciona uma dimensão ao
cálculo (mais granular) &
\texttt{\{INCLUDE [Product] : SUM([Sales])\}} —
total por produto, mesmo que a view mostre só categoria \\
\hline
\texttt{EXCLUDE} & Remove uma dimensão do
cálculo (menos granular) &
\texttt{\{EXCLUDE [Month] : AVG([Revenue])\}} —
média anual, mesmo que a view mostre mensal \\
\hline
\end{tabular}

\textbf{Quando usar:}
\begin{itemize}
  \item KPI de contexto que não deve mudar com filtros $\to$ \texttt{FIXED}
  \item Percentual do total por subgrupo $\to$ \texttt{FIXED} + cálculo de tabela
  \item Métrica por cliente independente da exibição $\to$ \texttt{FIXED [Customer]}
\end{itemize}
\end{BoardBox}

\begin{SolvedBox}
\textbf{Exemplo prático com World Happiness (professor demonstra):}

Calcular a média de felicidade por Region \textbf{fixa}, para usar como
benchmark mesmo quando o usuário filtra países individuais:

\texttt{\{FIXED [Region] : AVG([Happiness Score])\}}

Aplicação: adicionar como linha de referência no scatter PIB vs Felicidade.
Quando o usuário filtra para ``Latin America'', a linha de referência
continua mostrando a média da região completa, não apenas dos países
filtrados.
\end{SolvedBox}

% =============================================================================
\section{Dados Ausentes no Tableau (NULLs)}
% =============================================================================

\begin{BoardBox}
\textbf{Como o Tableau trata NULLs (para o quadro)}

Na Unidade I vimos dados ausentes (missing) no Python. No Tableau, os
mesmos problemas aparecem:

\begin{itemize}
  \item \textbf{NULLs em dimensões:} o Tableau cria uma categoria ``Null''
na visualização. Pode confundir ou esconder registros.
  \item \textbf{NULLs em medidas:} o Tableau ignora NULLs em
agregações (SUM, AVG). \texttt{AVG} de \{5, NULL, 10\} = 7.5
(não 5.0!). Isso pode ser enganoso se o NULL deveria ser 0.
  \item \textbf{NULLs em linhas temporais:} o Tableau pula datas
com NULL, criando ``buracos'' na linha. Solução:
Show Missing Values $\to$ marca as lacunas.
\end{itemize}

\textbf{Como investigar missingness no Tableau:}
\begin{enumerate}
  \item Campo calculado: \texttt{ISNULL([Campo])} retorna True/False
  \item Contar NULLs: \texttt{SUM(IF ISNULL([Campo]) THEN 1 ELSE 0 END)}
  \item Substituir NULL por valor:
  \texttt{IFNULL([Campo], 0)} ou \texttt{ZN([Campo])} (zero se null)
  \item \textbf{Visual de completude:} barras com \% de preenchimento
  por coluna — útil como primeira sheet de qualquer dashboard de AED
\end{enumerate}

\textbf{Regra de ouro:} documente a decisão. Se substituiu NULL por 0,
registre que fez isso e por quê. A audiência precisa saber.
\end{BoardBox}

% =============================================================================
\section{Parte IV — Ações de Dashboard (20h40–21h)}
% =============================================================================

\begin{BoardBox}
\textbf{Tipos de ações no Tableau (para o quadro)}

Ações transformam o dashboard em uma \textbf{experiência de exploração dirigida}.
O usuário clica em um elemento e o dashboard reage.

\begin{tabular}{|l|p{5.5cm}|p{4.5cm}|}
\hline
\textbf{Tipo} & \textbf{O que faz} & \textbf{Exemplo} \\
\hline
Filter Action & Clique num elemento filtra os outros gráficos & Clicar em ``SP'' no mapa filtra todos os KPIs por SP \\
Highlight Action & Destaca o elemento selecionado em todos os gráficos & Clicar em produto A ilumina todas as suas ocorrências \\
URL Action & Abre um link externo com parâmetros dinâmicos & Clicar num cliente abre o CRM naquele registro \\
Navigate Action & Navega para outro sheet/dashboard & Botão ``Ver detalhes'' vai para uma tela de drill-down \\
\hline
\end{tabular}

\textbf{Boas práticas de ações:}
\begin{itemize}
  \item Configure Run on: \texttt{Select} (mais intuitivo que \texttt{Hover}).
  \item Deixe claro visualmente que o elemento é clicável (tooltip ou ícone).
  \item Defina o comportamento ao desmarcar: ``Mostrar todos'' ou ``Excluir tudo''
  — nunca deixe o usuário preso num estado vazio sem saber como desfazer.
\end{itemize}
\end{BoardBox}

\begin{SolvedBox}
\textbf{Questão para o Quadro — Filtro, Parâmetro ou Ação?}

Uma varejista com 25 lojas quer um dashboard no Tableau que permita:
\begin{enumerate}
  \item Ver vendas apenas de uma região (Norte, Nordeste, Sul\ldots).
  \item Alternar a métrica principal entre ``receita'' e ``margem \%''.
  \item Ao clicar em uma loja no mapa, ver o detalhamento dela no
  gráfico de barras ao lado.
\end{enumerate}

Para cada necessidade, indique: \textbf{Filtro}, \textbf{Parâmetro}
ou \textbf{Ação de dashboard}?

\textbf{Resolução:}
\begin{center}
\begin{tabular}{|l|l|p{6.5cm}|}
\hline
\textbf{Necessidade} & \textbf{Recurso} & \textbf{Justificativa} \\
\hline
Selecionar região & Filtro & Subconjunto de dados (dropdown exposto) \\
\hline
Alternar métrica & Parâmetro & Controla lógica de cálculo, não dados \\
\hline
Clique no mapa & Ação & Filtragem interativa entre sheets \\
\hline
\end{tabular}
\end{center}

\textbf{Dica:} Filtro = ``quero ver menos dados''. Parâmetro =
``quero mudar \emph{o que} é calculado''. Ação = ``quero que clicar
aqui mude algo ali''.
\end{SolvedBox}

% =============================================================================
\section{Parte V — Construção Guiada: Dashboard Interativo com World Happiness (21h–21h40)}
% =============================================================================

\begin{SolvedBox}
\textbf{Construção guiada — ``Dashboard: O que Faz um País Feliz?''}

\textbf{Sheet 1 — Mapa Coroplético de Felicidade (10 min):}
\begin{enumerate}
  \item Arrastar \texttt{Country} para Detail (Tableau cria mapa automático)
  \item Arrastar \texttt{Happiness Score} para Color
  \item Paleta: sequencial (vermelho claro $\to$ verde escuro)
  \item \textbf{Título-mensagem:} ``Países nórdicos lideram o ranking de
  felicidade — países da África Subsaariana estão na base''
  \item Tooltip: país + score + GDP + Freedom
\end{enumerate}

\textbf{Sheet 2 — Scatter: PIB per Capita vs Felicidade (10 min):}
\begin{enumerate}
  \item Arrastar \texttt{GDP per Capita} para Columns,
  \texttt{Happiness Score} para Rows
  \item \texttt{Country} em Detail, \texttt{Region} em Color
  \item Analytics $\to$ Trend Line (linear)
  \item \textbf{Título:} ``Mais PIB = mais felicidade?
  A relação existe, mas satura nos países ricos''
  \item Anotar outliers: Costa Rica (feliz com PIB moderate),
  Qatar (rico mas não tão feliz)
\end{enumerate}

\textbf{Sheet 3 — Barras: Ranking Top 10 e Bottom 10 com Parâmetro (10 min):}
\begin{enumerate}
  \item Criar parâmetro \texttt{p\_métrica} (String, lista:
  Happiness Score, GDP, Social Support, Freedom, Generosity)
  \item Campo calculado \texttt{m\_din}:
  \texttt{CASE [p\_métrica] WHEN ``Happiness Score'' THEN [Happiness Score]
  WHEN ``GDP'' THEN [GDP per Capita] \ldots END}
  \item Barras horizontais por Country, ordenadas por \texttt{m\_din}
  \item Filtro Top 10 / Bottom 10 (usando Sets ou Index())
  \item O usuário alterna a métrica no parâmetro
\end{enumerate}

\textbf{Montar o Dashboard Interativo (10 min):}
\begin{enumerate}
  \item Mapa no topo (largura total)
  \item Scatter (esquerda) + Barras ranking (direita)
  \item Quick Filter por Region (dropdown)
  \item Filter Action: clique em região no mapa filtra scatter + barras
  \item Exibir controle do parâmetro \texttt{p\_métrica}
  \item Testar: clicar em ``Western Europe'' no mapa $\to$
  scatter mostra só países europeus
\end{enumerate}
\end{SolvedBox}

% =============================================================================
\section{Estudo Estatístico com World Happiness (21h40–21h50)}
% =============================================================================

\begin{BoardBox}
\textbf{Estudo Estatístico (para o quadro — discutir com os dados no Tableau)}

\begin{enumerate}
  \item \textbf{Correlação:} No scatter PIB vs Felicidade, qual o
  $R^2$ da trend line? A correlação é linear ou parece logarítmica?
  Troque para Social Support vs Felicidade — a correlação é mais forte?

  \item \textbf{Outliers:} Identifique 2 países que são outliers
  positivos (mais felizes do que o PIB sugere) e 2 negativos
  (menos felizes do que o PIB sugere). Por que isso importa para
  políticas públicas?

  \item \textbf{Comparação regional:} Use o filtro de Region.
  A média de felicidade da América Latina é mais alta que a da
  Ásia, apesar do PIB menor. Qual fator explica isso?
  (Dica: olhe Social Support e Freedom.)

  \item \textbf{Parâmetro de métrica:} Alterne entre Generosity
  e Corruption. Qual tem correlação mais forte com Happiness?
  Correlação não implica causalidade — mas sugere hipóteses.
\end{enumerate}
\end{BoardBox}

% ---------------------------------------------------------------------------
\section{Lançamento — Trabalho 6}
% ---------------------------------------------------------------------------

\begin{SolvedBox}
\textbf{Trabalho 6 — Dashboard Final (entrega: 20/05/2026):}
\begin{itemize}
  \item[$\square$] Dashboard no Tableau Public (link público obrigatório)
  \item[$\square$] Mínimo de 3 sheets integrados no dashboard
  \item[$\square$] Pelo menos 1 parâmetro e 1 filter action implementados
  \item[$\square$] Todos os títulos são mensagens (testados pelo Teste do Relance)
  \item[$\square$] Big Idea visível no dashboard (texto ou KPI principal)
  \item[$\square$] Pitch ao vivo de 3 minutos na Semana 12
\end{itemize}
\end{SolvedBox}

% =============================================================================
\section{Revisão Semanal — Questão tipo Prova II (Q3)}
% =============================================================================

\begin{NoteBox}
Hoje revisamos o conteúdo da \textbf{Questão 3} da Prova II:
filtros, parâmetros e ações de dashboard no Tableau.
\end{NoteBox}

\begin{SolvedBox}
\textbf{Cenário (estilo Q3 da Prova II):}

Uma rede de restaurantes tem 30 unidades em 5 estados do Nordeste.
O gerente de operações quer um dashboard que permita:
\begin{enumerate}
  \item Ver apenas as unidades de um estado específico.
  \item Alternar a métrica exibida entre ``ticket médio'' e ``satisfação''.
  \item Ao clicar em uma unidade no mapa, ver os detalhes dela em um
  gráfico adjacente.
\end{enumerate}

\textbf{Questão:} Para cada necessidade, indique se deve usar
\textbf{Filtro}, \textbf{Parâmetro} ou \textbf{Ação de dashboard}.
Justifique.

\textbf{Resolução:}
\begin{enumerate}
  \item \textbf{Filtro} — selecionar um subconjunto dos dados (estado)
  é o uso clássico de filtro. Pode ser exposto como dropdown no dashboard.
  \item \textbf{Parâmetro} — alternar entre métricas é uma decisão do
  \emph{usuário} que controla a lógica do cálculo (campo calculado que
  referencia o parâmetro). Não é um filtro porque não remove dados,
  apenas muda o que é exibido.
  \item \textbf{Ação de dashboard (filtro)} — clicar em um elemento
  visual para filtrar outro sheet é uma ação de filtro interativa.
  Configurada em Dashboard $\to$ Actions $\to$ Filter.
\end{enumerate}

\textbf{Resumo:}
\begin{center}
\begin{tabular}{|l|l|l|}
\hline
\textbf{Necessidade} & \textbf{Recurso} & \textbf{Por quê?} \\
\hline
Selecionar estado & Filtro & Subconjunto de dados \\
Alternar métrica & Parâmetro & Controla lógica, não dados \\
Clique interativo & Ação de dashboard & Filtragem entre sheets \\
\hline
\end{tabular}
\end{center}
\end{SolvedBox}

% ---------------------------------------------------------------------------
\section{Materiais Preparados}
% ---------------------------------------------------------------------------
\begin{itemize}
  \item Dataset World Happiness Report no AVA (link Kaggle também disponível)
  \item Tableau Workbook de referência (.twbx) com os 3 sheets base
  \item Slides: hierarquia de filtros animada; passo a passo de parâmetro
  \item Checklist de interatividade para o T6
\end{itemize}

% ---------------------------------------------------------------------------
\section{Próximo passo}
% ---------------------------------------------------------------------------

Na \textbf{Semana 11} trabalhamos \textbf{dashboards que forçam decisão}:
estrutura contexto→diagnóstico→recomendação, KPIs com targets, e o
\textbf{pitch de 3 minutos} — como apresentar sua análise para um executivo
em menos tempo do que um intervalo comercial.

Perguntem aos alunos: \textit{``Se o CEO da empresa que vocês estão analisando
tivesse 5 minutos para ver o dashboard de vocês, quais 3 perguntas ele faria?
O dashboard de vocês responde a essas 3 perguntas sem precisar de explicação?''}
